\chapter{Lười, trì hoãn và ảo tưởng cố gắng}
\markboth{Lười, trì hoãn và ảo tưởng cố gắng}{Lười, trì hoãn và ảo tưởng cố gắng}

\section{Tối nay em định học mà cầm điện thoại một chút thôi}
\begin{truyen}{Tối nay em định học mà cầm điện thoại một chút thôi}{Classroom Talk}
\chuhoa{H}{ọc sinh: ``Em mở điện thoại tính coi 5 phút nghỉ đầu óc thôi. Xong ngẩng lên là mất gần tiếng.''}
\textit{Student: ``I picked up my phone thinking I would rest my brain for just five minutes. Then I looked up and almost an hour was gone.''}\\[4pt]

\teacherpair{Vậy em bị thiếu thời gian hay thiếu thắng nổi cái tiện trước mắt?}{So are you short on time, or are you simply failing to beat the convenience right in front of you?}

\studentpair{Chắc cái thứ hai.}{Probably the second one.}

\teacherpair{Đúng. Nhiều khi mình không lười kiểu nằm im. Mình lười kiểu tự cho phép mình trượt dần.}{Exactly. Sometimes laziness is not lying still. Sometimes it is the laziness of quietly giving yourself permission to slide.}
\end{truyen}

\begin{baihoc}
Nhiều buổi học hỏng không vì thiếu quyết tâm lớn, mà vì thua mấy cái tiện nhỏ lặp đi lặp lại.
\textit{(Many study sessions fail not from a lack of grand determination, but from repeated surrender to small comforts.)}
\end{baihoc}

\begin{ghinhoanh}
\textbf{Short English to remember:}

Small distractions become big losses.

Delay grows when you keep permitting it.
\end{ghinhoanh}

\begin{tuhocnhanh}
1. Cái gì hay kéo em trượt khỏi bàn học nhất? \textit{(What most often pulls you away from study?)}

2. Em có cách nào chặn nó trước khi bắt đầu học không? \textit{(Can you block it before your study session starts?)}
\end{tuhocnhanh}
\ngancach

\section{Em làm bài muộn chứ không phải em không làm}
\begin{truyen}{Em làm bài muộn chứ không phải em không làm}{Classroom Talk}
\chuhoa{H}{ọc sinh: ``Em vẫn làm mà thầy. Chỉ là em hay làm sát giờ thôi, chứ em đâu bỏ.''}
\textit{Student: ``I still do the work, sir. I just do it right before the deadline. It is not like I abandon it.''}\\[4pt]

\teacherpair{Làm sát giờ rồi nộp thiếu, nộp ẩu, thức khuya vật vờ hôm sau, vậy khác bỏ bao nhiêu?}{If you do it at the last second, submit it half-finished, submit it carelessly, and then stumble through the next day from lack of sleep, how different is that from abandoning it?}

\studentpair{Nghe hơi đau.}{That hurts a little.}

\teacherpair{Đau vì đúng. Nhiều người không bỏ việc, chỉ bỏ chất lượng và bỏ sức khỏe trước.}{It hurts because it is true. Many people do not abandon the task itself. They abandon the quality and their health first.}
\end{truyen}

\begin{baihoc}
Trì hoãn lâu ngày làm mình tưởng mình vẫn ổn, cho tới khi kết quả và cơ thể cùng phản ứng.
\textit{(Procrastination can make you feel functional until both your results and your body start objecting.)}
\end{baihoc}

\begin{ghinhoanh}
\textbf{Short English to remember:}

Late work is not always real discipline.

Finishing badly is not the same as finishing well.
\end{ghinhoanh}

\begin{tuhocnhanh}
1. Em hay nộp đúng hạn hay hay tự cứu phút cuối? \textit{(Do you submit on time, or survive at the last second?)}

2. Trì hoãn đang làm hỏng điều gì ngoài điểm số? \textit{(What besides grades is procrastination damaging?)}
\end{tuhocnhanh}
\ngancach

\section{Em học theo cảm hứng nên mất nhịp suốt}
\begin{truyen}{Em học theo cảm hứng nên mất nhịp suốt}{Classroom Talk}
\chuhoa{H}{ọc sinh: ``Lúc có mood em học sung lắm. Nhưng không có hứng là em tắt luôn. Em sống kiểu đó riết nên lúc lên lúc xuống.''}
\textit{Student: ``When I am in the mood, I study really hard. But when I do not feel inspired, I shut down completely. That is why my rhythm keeps going up and down.''}\\[4pt]

\teacherpair{Nếu đợi hứng mới làm, sau này đi làm em cũng toang với đời.}{If you wait for motivation before doing important things, real life will hit you very hard later.}

\studentpair{Sao nặng vậy thầy?}{Why does that sound so harsh, sir?}

\teacherpair{Vì việc quan trọng hiếm khi chờ cảm hứng. Người bền là người làm được cả lúc không thích.}{Because important work rarely waits for motivation. Stable people are the ones who can still act even when they do not feel like it.}
\end{truyen}

\begin{baihoc}
Cảm hứng giúp bắt đầu nhanh. Kỷ luật mới giúp không bỏ dở giữa đường.
\textit{(Motivation helps you start. Discipline keeps you from collapsing halfway.)}
\end{baihoc}

\begin{ghinhoanh}
\textbf{Short English to remember:}

Motivation is helpful.

Discipline is dependable.
\end{ghinhoanh}

\begin{tuhocnhanh}
1. Lịch học của em có phụ thuộc tâm trạng quá nhiều không? \textit{(Does your study routine depend too much on mood?)}

2. Em có một thói quen nhỏ nào làm được dù không có hứng? \textit{(What small habit can you still do without motivation?)}
\end{tuhocnhanh}
\ngancach

\section{Em muốn giỏi nhanh, giỏi chậm nản lắm}
\begin{truyen}{Em muốn giỏi nhanh, giỏi chậm nản lắm}{Classroom Talk}
\chuhoa{H}{ọc sinh: ``Em ghét cảm giác tập hoài mà chưa hay. Em muốn thấy tiến bộ lẹ, chứ chậm quá em tụt mood.''}
\textit{Student: ``I hate practicing for so long and still not being good. I want to see fast improvement. If it is too slow, I lose motivation.''}\\[4pt]

\teacherpair{Em muốn giỏi thật hay muốn cảm giác mình giỏi nhanh?}{Do you want real skill, or do you want the feeling of becoming good quickly?}

\studentpair{Chắc em mê cái cảm giác hơn.}{Maybe I am more attached to the feeling.}

\teacherpair{Nhiều người bỏ cuộc không phải vì không tiến bộ, mà vì tiến bộ không kịp làm họ sướng.}{Many people do not quit because they are not improving. They quit because the improvement is not fast enough to satisfy their ego.}
\end{truyen}

\begin{baihoc}
Người bền là người chịu được giai đoạn chưa giỏi mà vẫn làm tiếp.
\textit{(A durable learner can tolerate being unskilled for a while and still continue.)}
\end{baihoc}

\begin{ghinhoanh}
\textbf{Short English to remember:}

Fast progress feels good.

Slow mastery lasts longer.
\end{ghinhoanh}

\begin{tuhocnhanh}
1. Em đang theo đuổi năng lực thật hay cảm giác thắng nhanh? \textit{(Are you chasing real ability or quick emotional reward?)}

2. Em có chấp nhận một quá trình chậm nhưng chắc không? \textit{(Can you accept a slow but solid process?)}
\end{tuhocnhanh}
\ngancach